% Метод условного градиента: y_k — вершина многогранника, минимизирующая
% линейную модель; шаг делается вдоль отрезка [x_k, y_k].
\begin{figure}[htbp]
\centering
\begin{tikzpicture}[scale=1.0, >={Stealth[length=2.2mm]}]

% допустимое множество
\fill[blue!8] (0,0) -- (4.6,-0.2) -- (5.4,1.8) -- (3.1,3.4) -- (-0.4,2.1) -- cycle;
\draw[blue!55!black, thick]
  (0,0) -- (4.6,-0.2) -- (5.4,1.8) -- (3.1,3.4) -- (-0.4,2.1) -- cycle;
\node[blue!55!black, font=\footnotesize] at (4.6,2.75) {$S$};

% линии уровня целевой функции
\foreach \s in {1.0, 0.62} {
  \draw[gray!55, rotate around={-25:(1.1,1.0)}]
    (1.1,1.0) ellipse ({1.9*\s} and {1.15*\s});
}
\fill[gray!55] (1.1,1.0) circle (1.2pt);

% текущая точка
\coordinate (xk) at (3.3,0.95);
\fill[black] (xk) circle (2pt);
\node[font=\footnotesize, anchor=north] at ($(xk)+(0.08,-0.1)$) {$x_k$};

% антиградиент
\draw[gray!70!black, thick, ->] (xk) -- ++(-1.0,0.05)
  node[above, font=\scriptsize, xshift=6pt] {$-\nabla\varphi(x_k)$};

% вершина y_k — минимум линейной модели (левая верхняя вершина)
\coordinate (yk) at (-0.4,2.1);
\fill[orange!75!black] (yk) circle (2pt);
\node[font=\footnotesize, anchor=east] at ($(yk)+(-0.1,0)$) {$y_k$};

% отрезок [x_k, y_k] и шаг
\draw[orange!75!black, very thick] (xk) -- (yk);
\coordinate (xk1) at ($(xk)!0.42!(yk)$);
\fill[orange!75!black] (xk1) circle (1.8pt);
\node[font=\footnotesize, anchor=south] at ($(xk1)+(0.1,0.1)$) {$x_{k+1}$};

\node[font=\scriptsize, orange!60!black, anchor=north] at (1.6,0.42)
  {$x_{k+1} = (1-\alpha_k)x_k + \alpha_k y_k$};

\end{tikzpicture}
\caption{Шаг метода условного градиента. Вспомогательная задача
$\nabla\T\varphi(x_k)x \to \min$ линейна, поэтому её решение $y_k$ — вершина
многогранника $S$. Направление $s_k = y_k - x_k$ допустимо по построению, а
новая точка берётся как выпуклая комбинация $x_k$ и $y_k$, что автоматически
оставляет её в $S$ при $\alpha_k \in [0,1]$.}
\label{fig:frank-wolfe}
\end{figure}
